% ==============================================================================
% SECCIÓN 5: COMUNICACIONES EN EL EQUIPO Y REUNIONES
% ==============================================================================
\section{Herramientas para Comunicaciones en el Equipo y Reuniones}

\subsection{Canales Oficiales del Equipo}
\begin{itemize}[leftmargin=*]
  \item \textbf{WhatsApp:} Principalmente utilizaremos WhatsApp para la organización y concreción de fechas y horarios, así como para la difusión de avisos inmediatos o urgencias imprevistas.
  \item \textbf{Discord:} Plataforma destinada a realizar videollamadas cuando no podamos hacer reuniones presenciales, compartir pantalla en sesiones técnicas y vertebrar discusiones temáticas.
  \item \textbf{Jira:} Utilizaremos la herramienta Jira para el seguimiento y control de las tareas en progreso, gestión del Product Backlog, diagramas de Gantt y órdenes del día de reuniones.
  \item \textbf{Correo Institucional (@correo.ugr.es y @go.ugr.es):} Canal oficial reservado para la comunicación formal con el profesorado y el personal directivo del centro.
\end{itemize}

\subsection{Protocolo de Reuniones y Elaboración de Actas}
\begin{itemize}[leftmargin=*]
  \item \textbf{Reuniones Semanales Internas:} Encuentro semanal del equipo (presencial o por Discord) convocado formalmente mediante Jira o WhatsApp.
  \item \textbf{Actas Formales:} Toda reunión oficial se documentará obligatoriamente según la plantilla estandarizada (\texttt{docs/plantillas/plantilla\_acta\_reunion.md}). El Catalogador y el Moderador redactarán el acta y la publicarán en \texttt{docs/reuniones/} con plazo aproximado para antes de la siguiente reunión.
\end{itemize}
